\documentclass{situationist-preprint}

\title{The Beetle in the Box}
\msid{LNG-07-b}
\msfamily{LNG-07}
\msdiscipline{Philosophy of Language}
\mstype{Companion manuscript}
\mskeywords{private language; Wittgenstein; beetle in the box; sensation language; rule-following}

\begin{document}
\maketitle

\begin{abstract}
In \S293 of the \emph{Philosophical Investigations} Wittgenstein imagines a community in which everyone has a box containing something called a \enquote{beetle}, which no one can see into but his own. LNG-07-a argues that a language whose rules only its speaker could know can be neither followed nor used to communicate. We read the beetle passage as a second route to the same conclusion, from the side of the object rather than of the rule. We state its argument as an invariance claim: varying the private contents of the boxes leaves the correct use of the word unchanged, so the contents form no part of what fixes its meaning. We then show that the passage does not deny the existence of private sensations, and answer the objection that each speaker can look into her own box.
\end{abstract}

\section{The passage}

Wittgenstein asks us to suppose that everyone has a box with something in it that they call a \enquote{beetle}; no one can look into anyone else's box, and each says he knows what a beetle is only by looking at his own. It would then be possible for everyone to have something different in the box, or for the thing to be constantly changing, or for the box to be empty. If the word nonetheless has a use in the people's language, that use is not the use of a name of a thing; the thing in the box has no place in the language game, not even as a something, and it drops out of consideration whatever it is \citep[§293]{Wittgenstein1953}. The target is the model on which a sensation word gets its meaning by naming an inner object that each speaker identifies in her own case \citep[§§243, 258]{Wittgenstein1953}.

\section{The invariance argument}

\begin{definition}\label{d:game}
A \emph{beetle community} is a community whose members use \enquote{beetle} in public: they say \enquote{there is a beetle in my box}, they ask one another about beetles, they correct one another's uses. A \emph{content assignment} specifies what, if anything, is in each box.
\end{definition}

\begin{assumption}[Inaccessibility]\label{a:inacc}
In a beetle community, whether a member's use of \enquote{beetle} is accepted or corrected by others does not depend on the content of any box.
\end{assumption}

The assumption is just the stipulation of the case: no one can look into another's box, so no one's acceptance or correction can depend on what is in it.

\begin{proposition}[Invariance]\label{p:inv}
Given Assumption~\ref{a:inacc} and Assumptions~2.3 and~2.4 of LNG-07-a, whether a member's application of \enquote{beetle} is correct is the same under every content assignment.
\end{proposition}

\begin{proof}[Argument]
Suppose the correctness of some application depended on the content of the speaker's box. Then the standard of correctness for her applications of the rule, to the extent that it differs across content assignments, would be fixed by a fact only she could know, since only she can look into her box. By Assumption~2.4 of LNG-07-a, if her applications have an independent standard, someone else could know what the rule requires; so either the standard is not independent, and by Assumption~2.3 she follows no rule, or what the rule requires does not turn on the content. In the second case correctness is fixed by what others could know, which by Assumption~\ref{a:inacc} does not vary with the content assignment.
\end{proof}

\begin{corollary}\label{c:drop}
If \enquote{beetle} has a use in a beetle community, the contents of the boxes form no part of what fixes its meaning.
\end{corollary}

\begin{proof}
If the word has a use, its applications are governed by a rule that members follow. By Proposition~\ref{p:inv}, what that rule requires does not vary across content assignments, including the assignment on which every box is empty. So the contents do not contribute to fixing the rule, and by Principle~P1 of LNG-01-a, on which meaning is fixed by use, they do not contribute to fixing the meaning.
\end{proof}

The corollary is Wittgenstein's remark that the thing in the box drops out. Its connection with LNG-07-a is direct. Theorem~3.1 there shows that no one follows a rule private to her. The beetle argument shows that, where there is a public rule, the private object plays no role in it. Together they exhaust the options for the object-and-designation model: either the rule is private, and cannot be followed; or it is public, and the object is idle.

\section{What the passage does not show}

It does not show that there are no beetles, or that sensations are nothing. Wittgenstein says that a sensation is not a something, but not a nothing either, and that the conclusion is only that a nothing would serve as well as a something about which nothing could be said \citep[§304]{Wittgenstein1953}. Corollary~\ref{c:drop} concerns what fixes meaning, not what a word refers to. A sensation word such as \enquote{pain}, whose use is tied to public criteria of injury, behavior and avowal, may refer to the speaker's sensation; what it cannot do is have its meaning fixed by the speaker's inner identification of that sensation. This is the restriction LNG-07-a draws when it says that the argument denies private rules for talking about sensations, not private sensations \citep{Kenny1973}.

Nor does it show that a solitary speaker cannot have a language. A solitary speaker whose rules are of a kind another could come to know follows rules, as Proposition~3.4 of LNG-07-a (Crusoe) shows. The beetle community is not solitary; its defect is that the purported referent is inaccessible, not that others are absent.

\begin{objection}
The analogy fails for sensations, because each of us can look into her own box. I know what pain is from my own case, and my word refers to what I find there.
\end{objection}

\begin{reply}
Looking into one's own box yields how things seem to the looker. If the meaning of the word were fixed by what each speaker finds, then whether a later application was correct would be settled by whether the speaker's present seeming matches her memory of an earlier one, and that is a further seeming, not an independent standard. This is the dilemma of Section~4 of LNG-07-a. The objection's other claim, that the word refers to what she finds, is granted: the beetle argument concerns meaning, not reference.
\end{reply}

\begin{objection}
Kripke reads the private language argument as a consequence of a sceptical solution on which rule-following consists in community agreement \citep{Kripke1982}. On that reading the beetle passage adds nothing.
\end{objection}

\begin{reply}
The argument given here does not depend on Kripke's reading. Assumptions~2.3 and~2.4 of LNG-07-a require that a rule's standard be independent of how things seem to the follower and knowable by another, not that it be constituted by agreement. The beetle passage adds a point Kripke's reading leaves implicit: even where a rule is public, the private item cannot enter into what it requires.
\end{reply}

\begin{objection}
Perhaps all the beetles are in fact alike, and the community's word names that common kind.
\end{objection}

\begin{reply}
If they are alike, the word applies to them; but by Proposition~\ref{p:inv} it would apply in the same way if they were not, and the community's practice would be unchanged. A fact that makes no difference to the correctness of any application is not a fact on which the meaning depends, however it happens to be.
\end{reply}

\section{Conclusion}

The beetle in the box shows that a private object cannot fix the meaning of a word with a public use, just as the argument of LNG-07-a shows that a private rule cannot be followed. A language could be private only by having one or the other, and so there is no private language. What survives is the private sensation, which public words can be about.

\begin{thebibliography}{9}\small
\bibitem[Kenny(1973)]{Kenny1973} Kenny, A. (1973). \emph{Wittgenstein}. Allen Lane.
\bibitem[Kripke(1982)]{Kripke1982} Kripke, S.~A. (1982). \emph{Wittgenstein on Rules and Private Language}. Harvard University Press.
\bibitem[Wittgenstein(1953)]{Wittgenstein1953} Wittgenstein, L. (1953). \emph{Philosophical Investigations}. Translated by G.~E.~M. Anscombe. Blackwell.
\end{thebibliography}

\end{document}
